\input{../tex-inputs/latex-leseansicht-vorspann}

\section[Arthur Schnitzler und Olga Gussmann an Gustav Schwarzkopf, 11. 6. 1903]{L04431 Arthur Schnitzler und Olga Gussmann an Gustav Schwarzkopf, 11. 6. 1903}
\nopagebreak\mylabel{L04431v}
\rehead{ }\normalsize\beginnumbering\briefempfaengerindex{Schwarzkopf, Gustav@\textsc{Schwarzkopf, Gustav}!zzzSchnitzler, Olga@\emph{von Olga Schnitzler}!1903-06-111@{11. 6. 1903}|(be}\briefempfaengerindex{Schwarzkopf, Gustav@\textsc{Schwarzkopf, Gustav}!zzzSchnitzler, Arthur@\emph{von Arthur Schnitzler}!1903-06-111@{11. 6. 1903}|(be}
\toendnotes[C]{\smallbreak\pagebreak[2]}
\correspDesc{Versand  durch Arthur Schnitzler, Olga Schnitzler am 11. 6. 1903 in Monte Generoso
\newline{}Übermittlung  am 11. 6. 1903 in Monte Generoso
\newline{}Zustellung  am 12. 6. 1903 in Wien
\newline{}Erhalt  durch Gustav Schwarzkopf am 12. 6. 1903 in Wien}\toendnotes[C]{\smallbreak}
\Standort{DLA, A:Schnitzler, HS.1985.1.1897.}
\physDesc{Bildpostkarte, 208 Zeichen
\newline{}Handschrift Arthur Schnitzler: Bleistift, deutsche Kurrent
\newline{}Handschrift Olga Schnitzler: Bleistift, deutsche Kurrent
\newline{}Versand: 1) Stempel: »\nobreak{}\oindex{Monte Generoso@\textbf{Monte Generoso}|pwk}Monte Generoso – Vetta, 11. VI. 03\nobreak{}«.   2) Stempel: »\nobreak{}\oindex{I., Innere Stadt@\textbf{I., Innere Stadt}|pwk}Wien 1{[}/1{]} 1, 13. \textcolor{gray}{6}. 0\textcolor{gray}{3}, 11½V–1N., Bestellt\nobreak{}«. }\toendnotes[C]{\smallbreak}\pstart{}{\pb}\textsc{Wien}\oindex{Wien@\textbf{Wien}|pw}.\pend{}\pstart{}Herrn Guſtav Schwarzkopf\pend{}\pstart{}Wien I\oindex{I., Innere Stadt@\textbf{I., Innere Stadt}|pw}.\pend{}\pstart{}Tiefer Graben 23\oindex{Wien@\textbf{Wien}!I., Innere Stadt@\textbf{I., Innere Stadt}!Tiefer Graben 23@\textbf{Tiefer Graben 23}|pw}.\pend{}{\bigskip}
\pstart
           \centering{}{\pb}\textcolor{gray}{\textbf{Monte Generoso\oindex{Monte Generoso@\textbf{Monte Generoso}|pw}. Vista sul M. Viso\oindex{Monte Viso@\textbf{Monte Viso}|pw} e Lago di
                     Varese\oindex{Lago di Varese@\textbf{Lago di Varese}|pw}}}\pend
           \vspace{1em}
\pstart
           \noindent{}{\pb}»\label{K_L04431-1v}\edtext{Das Maulthier ſucht im Nebel ſeinen
                  Weg\pwindex{Goethe, Johann Wolfgang von 28.\,8.\,1749 Frankfurt am Main – 22.\,3.\,1832 Weimar@\textsc{Goethe, Johann Wolfgang von} (28.\,8.\,1749 Frankfurt am Main – 22.\,3.\,1832 Weimar)!Kennst du das Land, wo die Zitronen blühn@\strich\emph{Kennst du das Land, wo die Zitronen blühn}|pwv}}{\lemma{\textnormal{\emph{Das … Weg}}}\Cendnote{\textnormal{Vers aus dem Gedicht \emph{Mignon}\pwindex{Goethe, Johann Wolfgang von 28.\,8.\,1749 Frankfurt am Main – 22.\,3.\,1832 Weimar@\textsc{Goethe, Johann Wolfgang von} (28.\,8.\,1749 Frankfurt am Main – 22.\,3.\,1832 Weimar)!Kennst du das Land, wo die Zitronen blühn@\strich\emph{Kennst du das Land, wo die Zitronen blühn}|pwk} (\emph{Kennst du das Land? wo die Citronen blühn}\pwindex{Goethe, Johann Wolfgang von 28.\,8.\,1749 Frankfurt am Main – 22.\,3.\,1832 Weimar@\textsc{Goethe, Johann Wolfgang von} (28.\,8.\,1749 Frankfurt am Main – 22.\,3.\,1832 Weimar)!Kennst du das Land, wo die Zitronen blühn@\strich\emph{Kennst du das Land, wo die Zitronen blühn}|pwk}) von Goethe\pwindex{Goethe, Johann Wolfgang von 28.\,8.\,1749 Frankfurt am Main – 22.\,3.\,1832 Weimar@\textsc{Goethe, Johann Wolfgang von} (28.\,8.\,1749 Frankfurt am Main – 22.\,3.\,1832 Weimar)|pwk}.}}}\label{K_L04431-1}« Und grüßt ſie herzlich. \spacefill\mbox{ArthSch}\pend
           
\pstart
           {[}hs. O. Schnitzler:{]} Die Luft daſelbſt wirkt lebhaft auf den Geiſt, Witz und Verſand.\pend
           \pstart Auf Wiederſehen! \spacefill\mbox{OS.}\pend{}\selectlanguage{ngerman}\endnumbering\briefempfaengerindex{Schwarzkopf, Gustav@\textsc{Schwarzkopf, Gustav}!zzzSchnitzler, Olga@\emph{von Olga Schnitzler}!1903-06-111@{11. 6. 1903}|)be}\briefempfaengerindex{Schwarzkopf, Gustav@\textsc{Schwarzkopf, Gustav}!zzzSchnitzler, Arthur@\emph{von Arthur Schnitzler}!1903-06-111@{11. 6. 1903}|)be}\mylabel{L04431h}
\begin{anhang}
\end{anhang}\newcommand{\dateiname}{L04431}\newcommand{\titel}{Arthur Schnitzler und Olga Gussmann an Gustav Schwarzkopf, 11. 6. 1903}\newcommand{\editorInnen}{Herausgegeben von Jahnke, SelmaMüller, Martin Anton}\input{../tex-inputs/latex-leseansicht-abspann}
